% ============================================================
%  SIMULACRO DE EXAMEN — TEMA 1 (PRIMER LAPSO)
%  Decimales, Operaciones Combinadas y Ecuaciones
%  Curso: 1er Año — CELAM
%  Autor: Ing. Gabriel Astudillo
%  Basado en guia_t1l1.tex (evaluación individual de 20 puntos)
%  Nota: SOLO el enunciado. No incluye clave de respuestas ni
%  indicadores de dificultad (sin estrellitas).
% ============================================================

\documentclass[12pt, letterpaper]{article}

% ── Paquetes ─────────────────────────────────────────────────

\usepackage{mathwizards-guia}
\mwguia{Simulacro de Examen — Tema 1}{Ing. Gabriel Astudillo $\cdot$ 1er Año}

% ── Comandos propios de este simulacro ───────────────────────
\newcommand{\N}{\mathbb{N}}

\begin{document}
% ============================================================

% ── Portada / Título ─────────────────────────────────────────
\mwportada{Simulacro de Examen — Tema 1}{Decimales, Operaciones Combinadas y Ecuaciones}{Evaluación individual $\cdot$ 20 puntos $\cdot$ 90 minutos}

\vspace{4pt}
\begin{cuadroinfo}
  \small
  \textbf{Nombre:} \rule{0.55\linewidth}{0.4pt} \quad \textbf{Fecha:} \rule{0.15\linewidth}{0.4pt} \\[6pt]
  \textbf{Tiempo disponible:} 90 minutos \quad$\bullet$\quad \textbf{Puntaje total:} 20 puntos \\[4pt]
  \textbf{Estructura:} 3 secciones $\cdot$ 10 ejercicios $\cdot$ 2 puntos cada uno. \\[4pt]
  \textbf{Instrucciones:}
  \begin{enumerate}[leftmargin=*]
    \item Lee cada ejercicio con calma antes de resolverlo.
    \item Muestra \emph{todos} tus pasos: se evalúa el procedimiento, no solo la respuesta final.
    \item No se permite el uso de calculadora.
    \item Cuida la \emph{posición} de la coma decimal y los \emph{ceros} no significativos.
    \item En las ecuaciones, comprueba tu solución cuando pertenezca a $\N$; si no pertenece, escribe $x \notin \N$.
    \item Escribe con letra clara y revisa tu examen antes de entregarlo.
  \end{enumerate}
\end{cuadroinfo}

\vspace{8pt}

% ============================================================
\section{Operaciones con Números Decimales (8 puntos)}
% ============================================================

\begin{enumerate}[leftmargin=*]

  \item \textbf{(2 puntos)} Calcula el valor exacto, alineando las comas y completando con ceros cuando sea necesario.
        \begin{enumerate}[label=\alph*), leftmargin=*]
          \item $1\,200\,000{,}0005 - 478\,653{,}09975$
          \item $65{,}408 + 2\,999{,}75 + 0{,}0009 + 147{,}00625$
        \end{enumerate}

  \item \textbf{(2 puntos)} Calcula el producto exacto, controlando el total de cifras decimales.
        \begin{enumerate}[label=\alph*), leftmargin=*]
          \item $0{,}0625 \times 4\,800$
          \item $3{,}75 \times 0{,}048$
          \item $0{,}000125 \times 6\,400$
        \end{enumerate}

  \item \textbf{(2 puntos)} Calcula el cociente, mostrando cómo mueves la coma del divisor.
        \begin{enumerate}[label=\alph*), leftmargin=*]
          \item $4{,}032 \div 0{,}064$
          \item $28\,125 \div 0{,}375$
          \item $85 \div 6$ \quad (cociente con dos cifras decimales por defecto; indica además el resto)
        \end{enumerate}

  \item \textbf{(2 puntos)} Calcula, respetando la jerarquía de las operaciones.
        \begin{enumerate}[label=\alph*), leftmargin=*]
          \item $0{,}125 \times 4\,000 \div 0{,}5$
          \item $2{,}4 \times 0{,}05 \div 0{,}006$
          \item $0{,}36 \div 0{,}08 \times 12{,}5$
        \end{enumerate}

\end{enumerate}

% ============================================================
\section{Operaciones Combinadas (6 puntos)}
% ============================================================

\begin{enumerate}[leftmargin=*, resume]

  \item \textbf{(2 puntos)} Resuelve respetando la jerarquía y los signos de agrupación.
        \begin{enumerate}[label=\alph*), leftmargin=*]
          \item $60 - \left\{\left[\,4^2 \times (9 - 5)\,\right] - \sqrt{196}\right\}$
          \item $\left\{\left[(3^4 - 6^2) \times 2\right] + \sqrt{400}\right\} \div 11$
        \end{enumerate}

  \item \textbf{(2 puntos)} Resuelve respetando la jerarquía y los signos de agrupación.
        \begin{enumerate}[label=\alph*), leftmargin=*]
          \item $80 - 3 \times \left\{\left[(12 - 7) \times 4\right] + \sqrt{81} \div 3\right\}$
          \item $7 \times \left\{\,6 + 4 \times \left[\,25 - 3 \times (11 - 8)\,\right]\right\}$
        \end{enumerate}

  \item \textbf{(2 puntos)} Resuelve respetando la jerarquía y los signos de agrupación.
        \begin{enumerate}[label=\alph*), leftmargin=*]
          \item $50 - 2 \times \left\{\left[(8{,}5 - 3{,}25) \times 4\right] \div 3 + 1{,}5\right\}$
          \item $0{,}6 \times \left[\,15 \div (2{,}5 - 1) + 2^3\,\right] - 1{,}4$
          \item $12{,}5 \times \left[(7 - 2)^2 \div 5\right] - 3{,}4$
        \end{enumerate}

\end{enumerate}

\newpage

% ============================================================
\section{Ecuaciones de Primer Grado (6 puntos)}
% ============================================================

\begin{enumerate}[leftmargin=*, resume]

  \item \textbf{(2 puntos)} Resuelve cada ecuación eliminando primero los signos de agrupación. Comprueba tu solución y, si no pertenece a $\N$, escribe $x \notin \N$.
        \begin{enumerate}[label=\alph*), leftmargin=*]
          \item $4(x - 3) + 2(x + 5) = 3(x + 7) - 5$
          \item $7x - 3(2x - 5) = 2(x + 4) + 9$
          \item $4(3x + 1) - 5(x - 2) = 6(x + 3) + 3$
        \end{enumerate}

  \item \textbf{(2 puntos)} Resuelve cada ecuación con fracciones. Comprueba tu solución y, si no pertenece a $\N$, escribe $x \notin \N$.
        \begin{enumerate}[label=\alph*), leftmargin=*]
          \item $\dfrac{x + 1}{2} + \dfrac{x + 2}{3} = 12$
          \item $\dfrac{3x - 2}{5} = \dfrac{x + 6}{2}$
          \item $\dfrac{4x - 3}{5} = \dfrac{x + 9}{2}$
        \end{enumerate}

  \item \textbf{(2 puntos)} Resuelve cada ecuación agrupando términos semejantes. Comprueba tu solución y, si no pertenece a $\N$, escribe $x \notin \N$.
        \begin{enumerate}[label=\alph*), leftmargin=*]
          \item $9x - 5 + 2x = 4x + 23 + 3x$
          \item $7x - 12 + 3x = 5x + 13$
          \item $9x - 7 + 2x = 3x + 25$
        \end{enumerate}

\end{enumerate}

\vspace{10pt}
\begin{cuadroinfo}
  \small
  \textbf{Recuerda:} revisa dos veces tus operaciones con decimales, respeta la jerarquía de las operaciones y, en las ecuaciones, sustituye el valor obtenido para verificar la igualdad. \textquestiondown Mucho éxito!
\end{cuadroinfo}

\end{document}
